\section{Algoritma \textit{Variational Quantum Eigensolver} (VQE)}

\subsection{Prinsip Variasi (Teorema Rayleigh-Ritz)}
Algoritma \textit{Variational Quantum Eigensolver} (VQE) beroperasi secara fundamental berdasarkan Prinsip Variasi atau Teorema Rayleigh-Ritz yang menetapkan landasan logis bagi pencarian energi terendah dalam sistem kuantum. Teorema ini menyatakan bahwa nilai ekspektasi energi dari sembarang fungsi gelombang percobaan (\textit{trial state}) $|\psi\rangle$ akan selalu bernilai lebih besar atau sama dengan energi keadaan dasar sejati ($E_0$) dari Hamiltonian sistem tersebut. Sebagaimana dirumuskan dalam persamaan (\ref{eq:energy_functional}), fungsional energi $E[\psi]$ dihitung sebagai rasio antara nilai ekspektasi Hamiltonian terhadap normalisasi keadaan sistem guna memastikan validitas fisik dari hasil pengukuran. Pencarian nilai minimum energi dilakukan melalui proses iteratif untuk mendekati batas bawah absolut tersebut, yang secara teoretis menjamin konvergensi menuju solusi optimal bagi masalah finansial yang telah dipetakan.
\begin{equation}
E[\psi] = \frac{\langle \psi | \hat{H} | \psi \rangle}{\langle \psi | \psi \rangle} \ge E_0
\label{eq:energy_functional}
\end{equation} (1)

Pembuktian formal terhadap Prinsip Variasi dilakukan dengan mengekspansi sembarang \textit{trial state} $|\psi\rangle$ ke dalam kombinasi linear dari basis keadaan eigen ortonormal $\{|n\rangle\}$ yang dimiliki oleh Hamiltonian $\hat{H}$. Melalui dekomposisi ini, nilai ekspektasi energi dapat dinyatakan sebagai jumlahan diskrit dari nilai eigen energi $E_n$ yang terbobot oleh kuadrat modulus amplitudo probabilitas masing-masing keadaan basis. Karena urutan tingkat energi sistem senantiasa memenuhi kondisi $E_0 \le E_1 \le E_2 \le \dots$, maka setiap komponen energi dalam ekspansi tersebut dapat digantikan dengan batas bawah absolut $E_0$ tanpa melanggar ketidaksamaan matematis. Hasil akhir dari derivasi ini mengonfirmasi bahwa nilai ekspektasi energi tidak akan pernah melampaui ke bawah ambang energi \textit{ground state}, sehingga memberikan jaminan integritas pada setiap solusi optimasi yang ditemukan.
\begin{equation}
\langle \psi | \hat{H} | \psi \rangle = \sum_n |c_n|^2 E_n \ge E_0 \sum_n |c_n|^2 = E_0
\label{eq:variational_inequality}
\end{equation} (2)

\subsection{Fondasi Ansatz: \textit{Unitary Coupled Cluster} (UCC)}
\textit{Unitary Coupled Cluster} (UCC) merupakan sebuah metodologi standardisasi dalam pembuatan fungsi gelombang percobaan yang terinspirasi secara mendalam oleh teori korelasi elektron di dalam bidang kimia kuantum. Metode ini menggunakan operator eksponensial uniter berbentuk $e^{T-T^\dagger}$ untuk mentransformasi keadaan referensi awal menjadi keadaan kuantum yang mengandung korelasi antar partikel secara eksplisit dan sistematis. Sifat uniter dari operator klaster ini sangat krusial karena menjamin bahwa sirkuit kuantum yang dihasilkan tetap bersifat reversibel serta mampu menjaga normalisasi probabilitas sepanjang seluruh proses optimasi parameter berlangsung. Penggunaan \textit{ansatz} yang berbasis pada prinsip fisik ini memberikan keunggulan dalam hal akurasi teoretis karena struktur sirkuitnya dirancang untuk menangkap interaksi fundamental yang ada di dalam Hamiltonian masalah.

Meskipun menawarkan akurasi teoretis yang tinggi, implementasi \textit{ansatz} UCC murni pada perangkat keras kuantum era \textit{Noisy Intermediate-Scale Quantum} (NISQ) menghadapi tantangan serius berupa kedalaman sirkuit yang sangat masif. Sirkuit yang terlalu dalam mengakibatkan akumulasi galat gerbang dan dekoherensi \textit{qubit} yang cepat, sehingga sering kali mendegradasi kualitas solusi akhir sebelum proses konvergensi parameter berhasil dicapai secara sempurna. Keterbatasan fisik ini memicu terjadinya transisi paradigma dari penggunaan \textit{ansatz} yang terinspirasi fisik murni menuju arsitektur yang dirancang khusus untuk efisiensi pada perangkat keras (\textit{hardware-efficient ansatz}). Jembatan logika ini sangat penting untuk dipahami guna menyeimbangkan antara kebutuhan ekspresivitas model dengan kapasitas operasional komputer kuantum yang tersedia saat ini di dalam ekosistem riset finansial.

\subsection{Arsitektur Ansatz \texttt{EfficientSU2}}
Arsitektur sirkuit parameterisasi \texttt{EfficientSU2} dirancang secara spesifik untuk mengoptimalkan efisiensi komputasi pada perangkat kuantum jangka pendek melalui pengurangan kedalaman sirkuit tanpa mengorbankan kapasitas eksplorasi ruang parameter. Sirkuit ini menggunakan kombinasi strategis dari gerbang rotasi $R_y(\theta)$ dan $R_z(\phi)$ yang merepresentasikan elemen grup uniter spesial $SU(2)$ pada setiap \textit{qubit} untuk membangun fleksibilitas orientasi vektor keadaan. Guna membangun korelasi antar variabel keputusan portofolio, lapisan keterbelitan yang menggunakan gerbang \textit{Controlled-NOT} ($CX$) diimplementasikan dengan pola konektivitas \textit{linear} yang menghubungkan \textit{qubit} bertetangga secara berurutan. Figure 13 menyajikan visualisasi arsitektur sirkuit \texttt{EfficientSU2} yang memperlihatkan struktur berlapis antara blok rotasi parameterisasi dan blok interaksi keterbelitan yang membentuk fungsi gelombang percobaan yang kompleks.

Fleksibilitas dari \textit{ansatz} \texttt{EfficientSU2} ditentukan oleh jumlah repetisi sirkuit ($d$) yang secara langsung mempengaruhi kemampuan model dalam mewakili solusi finansial yang memiliki dimensi tinggi. Semakin dalam kedalaman sirkuit yang digunakan, maka semakin besar pula kapasitas ekspresivitas \textit{ansatz} dalam menangkap hubungan kovariansi aset yang rumit, namun risiko akumulasi \textit{noise} sistemik juga akan meningkat secara proporsional. Struktur sirkuit ini memungkinkan pemetaan hubungan ketergantungan antar aset finansial ke dalam struktur keterbelitan fisik \textit{qubit}, sehingga menciptakan representasi sirkuit yang selaras dengan karakteristik masalah optimasi portofolio. Dengan mengatur parameter rotasi $\theta$ secara iteratif, sistem dapat menavigasi lanskap energi finansial guna menemukan konfigurasi alokasi kapital yang paling efisien berdasarkan data pasar yang diinputkan ke dalam Hamiltonian.

\subsection{Siklus Hibrida Klasik-Kuantum}
Algoritma VQE bekerja secara operasional sebagai sebuah sistem hibrida yang membagi beban tugas komputasi antara prosesor kuantum (QPU) dan prosesor klasik (CPU) guna mengatasi keterbatasan masing-masing platform. Unit prosesor kuantum bertanggung jawab penuh dalam melakukan preparasi keadaan percobaan $|\psi(\theta)\rangle$ pada sirkuit kuantum serta melakukan pengukuran nilai ekspektasi dari Hamiltonian biaya yang merepresentasikan profil risiko pasar. Hasil pengukuran dari QPU kemudian dikirimkan kembali ke unit prosesor klasik untuk diolah menggunakan algoritma optimasi stokastik guna menghitung gradien dan menentukan langkah pembaruan parameter $\theta$ selanjutnya. Pembagian tugas ini memungkinkan pemanfaatan keunggulan paralelisme kuantum dalam eksplorasi ruang keadaan sekaligus memanfaatkan stabilitas logika kontrol klasik dalam proses pengambilan keputusan optimasi.

Siklus hibrida ini terus berlangsung secara iteratif dalam sebuah \textit{feedback loop} yang kontinu hingga nilai energi sistem mencapai titik konvergensi minimum yang bersifat stabil dan memenuhi kriteria penghentian. Figure 14 mengilustrasikan diagram alir operasional dari loop hibrida VQE yang menunjukkan interaksi timbal balik antara domain kuantum dan domain klasik dalam proses pencarian solusi portofolio optimal. Konvergensi energi menuju nilai terendah menandakan bahwa sistem telah berhasil menemukan konfigurasi bitstring yang mewakili alokasi aset paling efisien sesuai dengan prinsip-prinsip optimasi Markowitz. Keberhasilan siklus ini membuktikan bahwa integrasi antara teknologi kuantum dan metode numerik klasik merupakan strategi yang paling robust untuk menyelesaikan tantangan optimasi finansial yang bersifat \textit{NP-hard} di masa depan.
